\documentclass{rptuma}
% Shared preamble for the three structure proposals. Built against ../rptuma.cls by build.sh,
% so each proposal renders exactly as a thesis written in the department template would.
\usepackage{xcolor}
\usepackage{booktabs}
\usepackage{tabularx}
\usepackage{enumitem}

% The class loads its logo as img/RPTUlogo; the image lives one directory up.
\graphicspath{{../}}

\definecolor{materialgray}{gray}{0.38}
\newcommand{\material}[1]{\par{\raggedright\small\sffamily\color{materialgray}\textbf{Material:} #1\par}}
\newcommand{\repo}[1]{\path{#1}}
\newcommand{\us}{\,\textmu s}

\setlist{nosep,leftmargin=*}
% Repository paths and identifiers are long unbreakable words; loose spacing beats margin overruns.
\AtBeginDocument{\sloppy}
\setcounter{tocdepth}{2}
% KOMA appends a period to section numbers once appendices exist; the department sample has none.
\KOMAoptions{numbers=noendperiod}
\addto\captionsenglish{\renewcommand{\bibname}{Sources Used for This Proposal}}

\newdate{thesisDate}{07}{10}{2026}
\author{Author's Name}
\department{University of Kaiserslautern-Landau\\Department of Computer Science\\Distributed Systems Lab}
\supervisortitle{Examiner:}
\supervisor{Prof. Dr.-Ing. Jens B. Schmitt}
\secondsupervisor{M. Sc. Supervisor Name}

\newcommand{\frontmatternote}{%
\section*{Front matter (unchanged from the department template)}
\begin{itemize}
  \item \textbf{Title page} and \textbf{Eigenst\"andigkeitserkl\"arung}, as generated by \texttt{rptuma.cls}.
        The declaration requires disclosing generative-AI use beyond grammar and style.
  \item \textbf{Abstract} (English): the problem (deception at the web edge has an unknown cost), the
        approach (one mechanism on four edges, levelled to identical behaviour), the method (paired
        benchmark on three measurement planes) and the headline numbers.
  \item \textbf{Zusammenfassung}: the same abstract in German.
  \item \textbf{Contents}, \textbf{List of Figures}, \textbf{List of Tables}.
  \item \textbf{Glossary} (optional, the class's \texttt{glossar} environment): WADM, honeytoken, edge,
        tier, VU, iteration, SLO, cgroup, LuaJIT, PUC Lua, proxy-wasm.
\end{itemize}
In the department template's own example, websites are cited through an archived copy
(\texttt{web.archive.org}); do the same for any web source in the list at the end.}


\subject{Master Thesis -- Structure Proposal C}
\title{Template C: Design Science -- Building a Fair Comparison}

\begin{document}

\maketitle
\pagenumbering{roman}
\setcounter{page}{1}

\chapter*{About This Proposal}

\textbf{Organising idea.} The thesis follows the design science research process: problem,
objectives, design and development, demonstration, evaluation \cite{peffers2007dsrm}. It presents
\emph{two} artefacts. Artefact~I is WADM, the deception layer. Artefact~II is the method that makes
four different edges comparable: behavioural and tuning levelling, plus a paired benchmark with
validity gates. The flaws found and fixed along the way become content, in the spirit of
benchmarking-flaw studies \cite{vanderkouwe2019sok}.

\textbf{Choose it if} you want the levelling and methodology work, a large share of the project's
effort (\texttt{EDGE\_LEVELING.md}, \texttt{BENCHMARK\_METHODOLOGY.md}), credited as a contribution
in its own right, and your examiners care most about measurement rigour.

\textbf{Trade-off.} This is the longest of the three structures, and it risks reading like a
chronicle. Keep the history to the flaws that changed a conclusion, each with before-and-after
evidence.

\begin{center}
\small
\begin{tabularx}{\textwidth}{@{}rXr@{}}
\toprule
Ch. & Title & approx.\ pages \\
\midrule
1 & Introduction & 5 \\
2 & Foundations & 9 \\
3 & Related Work & 6 \\
4 & Problem Analysis and Objectives & 7 \\
5 & Artefact I: The Deception Layer & 12 \\
6 & Artefact II: A Fair Cross-Edge Evaluation Method & 14 \\
7 & Demonstration & 6 \\
8 & Evaluation & 18 \\
9 & Discussion & 7 \\
10 & Conclusion and Future Work & 3 \\
A--D & Appendices & as needed \\
\bottomrule
\end{tabularx}
\end{center}
{\footnotesize Page counts are a rough budget, not a department rule.}

\textbf{How to read the rest.} The table of contents is the proposed structure. Under each
heading is a short note on what to write there, and a \emph{Material} line pointing to the
place in the repository that already holds it (\repo{results/} and \repo{plots/} are short for
\repo{benchmarks/results/} and \repo{benchmarks/results/plots/}).

\frontmatternote

\tableofcontents
\cleardoublepage
\pagenumbering{arabic}

%======================================================================
\chapter{Introduction}
\section{Problem and Motivation}
Deception at the edge needs no change to the application, but its cost is unknown, and comparing
it across edges is harder than it looks.

\section{Research Approach}
Design science with two artefacts, each built, demonstrated and evaluated in iterations.

\section{Contributions}
WADM on four edges; the levelling rules and parity check; the paired three-plane harness; the
measured costs; a catalogue of the comparison pitfalls found along the way.

\section{Outline}
One line per chapter, mapped to the design science steps.

%======================================================================
\chapter{Foundations}
\section{Deception and Honeytokens}
Honeypots, honeytokens, and why interaction with bait is a high-signal alert
\cite{spitzner2002honeypots,spitzner2003honeytokens}
\section{Edge Proxies and Extension Runtimes}
Request phases and buffering; LuaJIT, PUC Lua and WebAssembly.
\section{Performance Evaluation Foundations}
Latency and its percentiles, cgroup CPU accounting, open- and closed-loop load, pairing,
replication, crossover designs \cite{jain1991art,williams1949designs}.

%======================================================================
\chapter{Related Work}
\section{Deception at the Application Layer}
Decoys, honeywords and deceptive web artefacts, and deploying them without developer involvement
\cite{bowen2009baiting,juels2013honeywords,virvilis2014changing,kahlhofer2024application}
\section{Measuring Proxy Overhead}
How sidecar and proxy costs have been measured, and which of their choices this harness adopts
\cite{zhu2022dissecting}.
\section{Benchmarking Flaws in Systems Security}
\cite{vanderkouwe2019sok,heiser_crimes}. Which of the documented flaws this project ran into
itself.

%======================================================================
\chapter{Problem Analysis and Objectives}
Design science step 2: the objectives against which the artefacts are later judged.

\section{Scenario and Stakeholders}
An operator protects a web application whose code it cannot change.

\section{Threat Model}
A reconnaissance client reads HTML, replays values, tampers with cookies and form fields, and
probes for SQL injection.

\section{Functional Objectives}
F1 plant bait; F2 detect interaction with it; F3 strip the trigger before it reaches the origin;
F4 raise one alert per detection; F5 driven only by configuration.

\section{Non-Functional Objectives}
N1 low added latency; N2 low CPU cost; N3 no broken responses; N4 portable to four edges.

\section{Comparability Objectives}
C1 identical behaviour; C2 identical timed regions; C3 levelled product defaults; C4 a baseline
that differs only by WADM.

%======================================================================
\chapter{Artefact I: The Deception Layer}
\section{Architecture and Configuration}
Edge between client and static origin; one shared \texttt{config.json}.
\material{\repo{CONTEXT.md}, \repo{config.json}}
\section{Honeytoken Kinds}
\subsection{HTML Comments and HTTP Headers}
Planted secrets, detected when the keyword is replayed.
\subsection{Cookies and Form Fields}
Planted values, detected when tampered with.
\subsection{Decoy Paths}
An advertised link to a trap path; any request to it matches.
\section{The Fake SQL-Injection Trap}
A detection-only feature: the edge answers itself with a fake MySQL error.
\section{Realisation on Four Edges}
\subsection{OpenResty}
The reference edge: one precompiled table, four phase hooks, a shared dictionary for IPs.
\material{\repo{nginx/README.md}}
\subsection{Envoy with Lua}
One filter script; the request path is handed to the response phase.
\material{\repo{envoy_scripts/README.md}}
\subsection{Envoy with WebAssembly}
Rust proxy-wasm; fewer host calls; the trap still contacts the origin (a host constraint).
\material{\repo{wasm-filter/README.md}}
\subsection{Apache httpd with mod\_lua}
A shared core for three hooks; per-thread Lua states; PUC Lua rather than LuaJIT.
\material{\repo{apache_scripts/README.md}}

%======================================================================
\chapter{Artefact II: A Fair Cross-Edge Evaluation Method}
The core methodological contribution.

\section{Levelling Behaviour}
Differences found: Apache never stripped the keyword (Lua pattern characters in the keyword),
four different detection mechanisms, attacker IPs recorded as \texttt{unknown} on Envoy, a log
line on every request, PCRE against plain find. The canonical mechanism that replaced them.
\material{\repo{docs/THESIS_NOTES.md} (2026-09-18), \repo{docs/EDGE_LEVELING.md}}

\section{Levelling Product Defaults}
Upstream keep-alive (about 0.5\,ms per request on OpenResty), access logs, \texttt{LuaCodeCache}.
The rule this establishes: level defaults before comparing.
\material{\repo{docs/THESIS_NOTES.md} (2026-09-17)}

\section{Baselines and Tiers}
Bare (WADM deleted, nothing added), silent (as deployed), instrumented (timers on).

\section{Measurement Planes}
Timers, end-to-end latency, cgroup CPU, and why no single plane survives a noisy VM.

\section{From a Flawed Harness to a Paired One}
\subsection{Flaws Found}
Request position read as WADM phase, a mean dominated by its tail, unpaired tiers, unequal offered
load, environment variables silently overriding the scenario, log writes inside the measured path,
the load generator competing with the edge.
\subsection{Fixes}
Pairing, balanced tier and probe orders, medians and trimmed means, separate tiers for
instrumentation.
\material{\repo{docs/BENCHMARK_METHODOLOGY.md} \S1, \S7}

\section{Validity Gates and Negative Controls}
Gates scale with the damage rather than firing on any failure; origin CPU is the negative control.
\material{\repo{benchmarks/wadm_timings.py}}

%======================================================================
\chapter{Demonstration}
Design science step 4: the artefacts work as intended.

\section{Behavioural Parity}
\texttt{parity\_check.py}: identical responses, alerts and origin requests on all four edges.
\material{\repo{benchmarks/parity_check.py}}

\section{Functional Behaviour across Response Shapes}
Thirteen response types, bare against WADM: preserved, injected, changed or broken.
\material{\repo{results/functional_matrix.md}}

\section{An Attack Walk-Through}
Read the HTML comment, replay its keyword, tamper with a cookie, try SQL injection; the alerts
that result on each edge.

%======================================================================
\chapter{Evaluation}
Design science step 5: Artefact~I judged against N1--N3, Artefact~II against C1--C4.

\section{Setup and Protocol}
Testbed, workload, rates, replicates, rejected cells.
\material{\repo{docs/BENCHMARK_METHODOLOGY.md} \S4, \S7}
\section{Latency}
$+109$ to $+218$\us{} per \texttt{GET} at the median.
\material{\repo{plots/paired_latency_added.png}}
\section{CPU, Cost Attribution and the Probe Effect}
$+33$ to $+200$\us{} of CPU per request; the timers enclose 5--22\% of it.
\material{\repo{plots/cost_attribution.png}}
\section{Capacity}
Capacity kept with silent WADM, and the resource that bound it.
\material{\repo{results/capacity_summary.md}}
\section{Mechanisms, Page Size and the SQLi Trap}
The \texttt{GET} path against the trap, linear growth with body size, the 2$\times$2.
\material{\repo{results/mechanisms.md}, \repo{results/sqli.md}}
\section{Stability: Apache mod\_lua}
Segfaults driven by cumulative requests, and the \texttt{LuaScope conn} trade-off.
\material{\repo{plots/apache_remediation.png}}
\section{Evaluating the Method Itself}
How much levelling and pairing changed the numbers. WASM's 500-VU p90 fell from 382 to 65\,ms once
behaviour was levelled. The old harness could not resolve the effect it claimed to measure, while
the paired harness resolves every \texttt{GET} cell. The two-tier figures had overstated WADM's
cost by the instrumentation's own cost.
\material{\repo{docs/BENCHMARK_METHODOLOGY.md} intro and \S8.2}

%======================================================================
\chapter{Discussion}
\section{Meeting the Objectives}
A table mapping each of F1--F5, N1--N4 and C1--C4 to its evidence, with a verdict.
\section{Lessons for Cross-Platform Benchmarks}
Level the defaults, pair the tiers, gate on damage, keep a negative control.
\section{Threats to Validity and Limitations}
Single VM host, synthetic traffic, POST-body inspection on only two edges, no TLS.

%======================================================================
\chapter{Conclusion and Future Work}
\section{Summary}
What each artefact achieved, against the objectives.
\section{Future Work}
Complete the fourth edge's POST inspection, injection into compressed bodies, a field study, and
reusing Artefact~II for other edge features.

%======================================================================
\appendix
\chapter{Configuration Reference}
Every \texttt{config.json} field per honeytoken kind, and the SQLi trap's fields.
\material{\repo{CONTEXT.md}}
\chapter{Equivalence Invariants Checklist}
Each invariant, how it is checked, and the residual differences left on purpose.
\material{\repo{benchmarks/README.md} (``Level-playing-field invariants'')}
\chapter{Harness History and Archived Results}
Each archived result set and why it was superseded.
\material{\repo{benchmarks/results/archive/}}
\chapter{Reproduction Guide}
Commands and machine configuration.
\material{\repo{docs/BENCHMARK_METHODOLOGY.md} \S3--4}

\bibliographystyle{ieeetr}
\bibliography{sources}
\end{document}
